\section{Introduction}

Machine learning classifiers can change their prediction under small input perturbations, raising concerns about safety, security, and ethics. 
A central notion for studying this phenomenon is \emph{global robustness}:
 a classifier is $\epsilon$-robust if it assigns the same label to any two inputs within
distance~$\epsilon$~\cite{casadio2022neural}. 
The distance encodes similarity and differs by domain. 
In computer vision, distances on pixel intensities are common~\cite{wang2026survey}. 
In fairness, the distance encodes context-dependent similarity of individuals, for instance, being equally qualified for a job~\cite{fairness-landscape,kusner2017counterfactual,ranking-fairness}.
In natural language processing, it encodes syntactic and semantic similarity under meaningful linguistic variation~\cite{levy2025statistical}.
Importantly, these notions of similarity may \emph{not} be metric, as they may violate the triangle inequality~\cite{garcia2019learning,similarity-guide}.
We therefore study robustness over arbitrary \emph{semimetric} spaces.
On connected domains such as $\mathbb{R}^n$, every $\epsilon$-robust classifier
is essentially constant: overlapping $\epsilon$-balls form chains between any two inputs,
and the label must agree along each chain.
Existing work has addressed this limitation by weakening the notion: 
local robustness restricts it to hand-picked
inputs~\cite{li2024certifying}; confidence-based notions to high-confidence
inputs~\cite{athavale2024verifying,leino2021globally}; probabilistic notions bound the probability of
drawing an $\epsilon$-close pair with different labels~\cite{mangal2019robustness,levy2023groma}.
Probabilistic bounds, however, constrain  the probability of robustness violations, not their geometric structure: Fig.~\ref{fig:grid} shows a classifier that satisfies such a
bound although no input is locally robust.

To address this limitation, we introduce $(\mu,\epsilon)$-robustness:
Every input must lie within distance $\mu$ of a point with the same label at which the classifier is locally $\epsilon$-robust. 
The property recovers $\epsilon$-robustness at $\mu = 0$.
Unlike probabilistic relaxations, it constrains where non-robust points may lie by imposing an explicit coverage structure on non-robust points:
they form bands of width at most~$\mu$ around locally robust points.
Two of our main results, a covering characterization of robustness and a necessary-and-sufficient criterion for the
existence of a non-constant robust classifier, hold over arbitrary semimetric spaces,
providing a domain-independent foundation.
Further, the impossibility result above is a consequence of \emph{connected
continuous} domains, whereas deployed classifiers compute over finite-precision
arithmetic. In fact, verification methods reasoning over reals can be unsound 
due to their abstraction~\cite{jia2021exploiting,cordeiro2025neural} showing that the fact that deployed systems operate over finite representable values rather than continuous domains has significant implications.
On a finite domain, $\epsilon$-balls can only build chains if $\epsilon$ exceeds the gaps between discrete values, so
whether a non-constant robust classifier exists becomes a quantitative question
about $\epsilon$, $\mu$, and the spacing between values of the domain.

\paragraph{Contributions.}
Studying $(\mu,\epsilon)$-robustness gives us domain-independent structural results and quantitative guarantees that can be specialized (see Section~\ref{sec:modeling-arithmetic}) to finite-precision arithmetic:
\begin{enumerate}
\item \textbf{Covering characterization.}
We first introduce $(\mu,\epsilon)$-robustness (Section~\ref{sec:new-robustness})
and then present a covering characterization of $(\mu,\epsilon)$-robustness (Section~\ref{sec:partitions-covering}), formally proving and extending the classical connection between $\epsilon$-robustness and separability~\cite{yang2020closer,mangal2023certifying}.

\item \textbf{Existence criteria.}
We first introduce a floating-point model (Section~\ref{sec:modeling-arithmetic}), and
then 
derive necessary and sufficient conditions for the existence of a non-constant robust classifier (Section~\ref{sec:existence}), both for arbitrary semimetric spaces and for finite totally ordered domains.

\item \textbf{Precision-dependent bounds.}
We obtain bounds on $\epsilon$ and $\mu$ in terms of the largest and smallest representable gaps (Section~\ref{sec:upper-bounds}). These bounds explain empirically observed robustness gains under coarser quantization~\cite{xiao2023robustmq,shrestha2024properties,duncan2020relative}, and extend to some partially ordered domains.

\item \textbf{Mechanization.} All results are mechanized in Lean~4\footnote{\url{https://github.com/verinncoq/mueps-robustness-proofs},
  archived in LAX as \url{https://laxarchive.org/lax-544225/paper.html}.}. 
  The machine-checked proofs were
  developed with AI assistance, and we manually validated
  that the formal statements match those in the document. To connect our
  floating-point model (Section~\ref{sec:modeling-arithmetic}) to an established IEEE~754 semantics, an AI assistant autonomously
  verified the constructions in Rocq on top of
  Flocq~\cite{boldo2011flocq}, and differentially tested the resulting executable
  distances against an independent implementation on special values and
  randomly generated binary 64-bit patterns.

\end{enumerate}
Our results concern whether robust classifiers \emph{exist} in a given domain;
whether an \emph{accurate} robust classifier exists and how robust
models are trained and verified are orthogonal to our focus.